\paragraph{Exercise 1.} \(\varphi(S_+)\not\subseteq\mathfrak q\) means that some homogeneous \(f\in S_+\) satisfies \(\varphi(f)\notin\mathfrak q\), hence \(\mathfrak q\in D_+(\varphi(f))\).

\paragraph{Exercise 2.} A point \(\mathfrak q\) is bad exactly when \(u\in\mathfrak q\).  In \(\Proj k[u,v]\), this is the point \((u)\); every point of \(D_+(u)\) is good.

\paragraph{Exercise 3.} If a projective prime contained \(\varphi(S_+)T\), it would contain its radical and hence \(T_+\), contradiction.

\paragraph{Exercise 4.} Define \(T^{(d)}_r=T_{dr}\).  Then \(\varphi(S_r)\subseteq T^{(d)}_r\), so degrees are preserved.

\paragraph{Exercise 5.} The induced map is defined where not all images of the target variables lie in the projective prime, exactly the complement of the common vanishing locus.

\paragraph{Exercise 6.} The primes of \(B=A/I\) correspond to primes of \(A\) containing \(I\), giving a homeomorphism onto \(V(I)\); the sheaf map is locally the quotient map.

\paragraph{Exercise 7.} Localization gives \(S_f\twoheadrightarrow(S/I)_{\bar f}\) with kernel \(IS_f\).  Taking degree zero yields the stated quotient.

\paragraph{Exercise 8.} The quotient ring is \(k[x_0,\ldots,x_{n-1}]\), so its Proj is \(\mathbb P^{n-1}\) and the image is \(x_n=0\).

\paragraph{Exercise 9.} On affine charts, base changing \(A\twoheadrightarrow A/I\) along \(A\to A'\) gives \(A'\twoheadrightarrow A'/IA'\), again a quotient.

\paragraph{Exercise 10.} If \(X\hookrightarrow\mathbb P^n_Y\), then after base change \(X_{Y'}\hookrightarrow\mathbb P^n_{Y'}\) remains closed.

\paragraph{Exercise 11.} The map has kernel \((Y_0Y_2-Y_1^2)\); compare dimensions or reduce monomials to see there are no further relations.

\paragraph{Exercise 12.} Repeatedly choose at most \(d\) units from the exponent multiset to create each factor; after \(r\) steps all exponents are exhausted.

\paragraph{Exercise 13.} \(N+1=\binom{4}{2}=6\), with coordinates \(x_0^2,x_0x_1,x_0x_2,x_1^2,x_1x_2,x_2^2\).

\paragraph{Exercise 14.} The degree-zero localization is generated by \(z_{aj}/z_{ij}=x_a/x_i\) and \(z_{ib}/z_{ij}=y_b/y_j\), giving the tensor product of the two affine chart rings.

\paragraph{Exercise 15.} Under the map, the minor becomes \(x_i y_j x_\ell y_r-x_i y_r x_\ell y_j=0\).

\paragraph{Exercise 16.} Minor relations interchange two column indices while preserving row indices.  Thus monomials with identical row/column multisets are congruent; grouping a kernel polynomial by image monomials proves generation.

\paragraph{Exercise 17.} The coordinate ring is the quotient of \(k[Z_{00},Z_{01},Z_{10},Z_{11}]\) by the sole minor \(Z_{00}Z_{11}-Z_{01}Z_{10}\).

\paragraph{Exercise 18.} The polynomial relations defining Segre are universal over \(\mathbb Z\); tensoring with the base scheme gives the relative closed immersion.

\paragraph{Exercise 19.} Use \(X\hookrightarrow\mathbb P^m_Y\hookrightarrow\mathbb P^m_Z\times_Z\mathbb P^n_Z\hookrightarrow\mathbb P^N_Z\), where the middle inclusion is obtained from \(Y\hookrightarrow\mathbb P^n_Z\) by base change.

\paragraph{Exercise 20.} \(\mathbb P^0_Y\cong Y\), so the given closed immersion already provides the required projective factorization.

\paragraph{Exercise 21.} On a standard chart, a Veronese coordinate is a degree-\(d\) monomial.  Ratios of such coordinates transform by the \(d\)-th power of the transition factor for \(\OO(1)\), which is the transition law for \(\OO(d)\).

\paragraph{Exercise 22.} A matrix in \(\mathrm{GL}_{n+1}(k)\) acts invertibly on the span of the variables.  Extension to the symmetric algebra gives a graded automorphism; scalar matrices act trivially on projective points, leaving \(\mathrm{PGL}\).

\paragraph{Exercise 23.} For example \([x_0:x_1:x_2]\dashrightarrow[x_0x_1:x_1x_2]\) has base locus \(V_+(x_1)\cup V_+(x_0,x_2)\); on its complement it is a morphism.

\paragraph{Exercise 24.} Two schemes may have the same \(k\)-points but different nilpotent structure.  A closed immersion requires local surjections of structure rings, information invisible to point-set injectivity.
